\chapter[Python Code]{Python Code} \label{code}
\definecolor{mygreen}{rgb}{0,0.6,0}
\definecolor{mygray}{rgb}{0.5,0.5,0.5}
\definecolor{mymauve}{rgb}{0.58,0,0.82}

\lstset{ %
    backgroundcolor=\color{white},   % choose the background color
    basicstyle=\ttfamily\small,        % size of fonts used for the code
    breaklines=true,                 % automatic line breaking only at whitespace
    captionpos=b,                    % sets the caption-position to bottom
    commentstyle=\color{mygreen},    % comment style
    escapeinside={\%*}{*)},          % if you want to add LaTeX within your code
    keywordstyle=\color{blue},       % keyword style
    stringstyle=\color{mymauve},     % string literal style
    numberstyle=\tiny,
    breaklines=true,
    postbreak=\mbox{\textcolor{red}{$\hookrightarrow$}\space},
    language=Python,
    frame=single,
    columns=fullflexible,
    numbers=left,
}
\begin{lstlisting}[
]
# pre-existing model parameters
tau = 0.92
a1 = 0.98
a4 = 0.8
d2 = 0.11

d3 = 0.4
d1 = 0.29
c1 = 0.9 #9e-1
c3 = 0.9 #9e-1

c2 = 85e-3
c4 = 85e-3
rho = 0.1
alpha =  0.2

k=0.036
n = 3
lambda1 = 126.12
lambda2 = 0.85

k1 = 0.47
k2 = 0.57
k3 = 0.49
k4 = 0.061

r1 = 0.73
r2 = 0.27

# custom parameters
aq = 0.00125 # production term for TM->TQ
dq = 0.00125 # exit term for TQ->TI
tau_q = 30 

u_target = 0.5
u_active = 3
u_rest = 4
u_cycles = 8
u_start = 90

u_cycle_duration = u_active + u_rest # do not change

def delay_rhs(x, t):
    global dq
    TI, TM, I, u, TQ = x(t)
    TI_past, _, _, _, _ = x(t - tau)
    _, _, _, _, TQ_past= x(t - tau_q)


    dx = [0,0,0,0,0]
    dx[0] = 2*a4*TM - (c1*I + d2)*TI - a1*(TI_past) + dq*(TQ_past)# TI
    dx[1] = a1*(TI_past) - d3*TM - a4*TM - c3*TM*I - k1*(1 - np.exp(-k2*u))*TM - aq*TM # TM
    
    if (alpha * 1000 < TI + TM):
        dx[2] = k + rho * I - c2*I*TI - c4*TM*I - d1*I - k3*(1 - np.exp(-k4 * u))*I# I
    else:
        dx[2] = k + (rho * I * (TI + TM)**n)/(alpha + (TI + TM)**n) - c2*I*TI - c4*TM*I - d1*I - k3*(1 - np.exp(-k4 * u))*I# I

    dx[3] = - lambda2 * u # u
    if ( t > u_start and t < u_start + u_cycle_duration * u_cycles and (t - u_start) % u_cycle_duration < u_active):
        dx[3] = u_target
    
    dx[4] = 2*aq*TM - dq*(TQ_past) # TQ

    for i in range(4):
        if x(t)[i] + dx[i] >= 100 :
            dx[i] = 0

    return dx

# Initial history
def initial_history_func(t):
    #         TI TM I u TQ
    return [2.5, 2.5, 0.9, 0, 0]

# Time points we want
t = np.linspace(0, 2000, 40000)

# Perform the integration
x = ddeint(delay_rhs, initial_history_func, t)

plt.figure()
plt.plot(t, x[:,3], label="U (meds)", color='#4eb3e4')
plt.plot(t, x[:,0], label="TI (interphase)", color='#fabb4b')
plt.plot(t, x[:,1], label="TM (mitosis)", color='#e6222e')
plt.plot(t, x[:,2], label="I (immune)", color='#1d9a4d')
plt.plot(t, x[:,4], label="TQ (quiescent)", color='purple')
plt.plot(t, x[:,0]+ x[:,1]+ x[:,4], label="Total tumor", color='black')

plt.xlabel("Time")
plt.ylabel("Concentration")
plt.title(f"rest period = {u_rest}")
plt.legend()
plt.ylim(0,0.5)
plt.show()
\end{lstlisting}